% Provenance (never printed): NSI premiere, "Langages et programmation" (specification, preconditions and
% postconditions, assertions, test sets, a passing test set does not prove a program correct); NSI premiere,
% "Valeurs booleennes" (sequential evaluation of and/or); NSI terminale, "Mise au point des programmes"
% (typical causes of bugs). See sources/book2/BRIEFS.md.
\chapter{Specifications, Assertions and Tests}\label{ch:g10:specifications-and-tests}

Sam writes a function that tells whether a year is a leap year, the kind of year with a 29th of February. She
tries it on $2024$ (leap), on $2023$ (not leap) and on $2000$ (leap): three right answers. She shares it with her
class. A week later a friend who is building a calendar of the twentieth century finds that her function
gives $1900$ a 29th of February, a day that never existed.

Was the function wrong? It depends on what it promised. This chapter is about writing down that promise, the
\emph{\omterm{def:g10:specifications-and-tests:specification}{specification}}, so that ``right'' and ``wrong'' have a meaning; about making the program check the
promise while it runs, with \emph{\omterm{def:g10:specifications-and-tests:assertion}{assertions}}; and about choosing the inputs that are most likely to break a
function, the \emph{tests}, and running them automatically every time the code changes.

\begin{recall}
A function has a name, parameters and a body; \code{return} sends a value back to the caller (Book 1, grade 7).
A program can fail in three ways: it does not start, it stops with an error message, or it gives a wrong answer;
a trace table follows its variables step by step (Book 1, grade 7). A comparison such as \code{x < 10} has a
Boolean value, \code{True} or \code{False}, and \code{and}, \code{or}, \code{not} combine such values (Book 1,
grade 6).
\end{recall}

\section{What a function promises}\label{sec:g10:specifications-and-tests:promise}

A function is written once and called many times, often by people who never read its body. They need to know two
things: what they must give it, and what they get back.

\begin{definition}[Specification]\label{def:g10:specifications-and-tests:specification}
The \emph{specification}\index{specification} of a function states, without describing how the function works,
what it computes: its name and parameters, the values its parameters may take, and the result it returns or the
change it makes for each of those values.
\end{definition}

A \omterm{def:g10:specifications-and-tests:specification}{specification} is usually split into two parts: what the caller must guarantee, and what the function
guarantees in return.

\begin{definition}[Precondition, postcondition]\label{def:g10:specifications-and-tests:pre-post}
A \emph{precondition}\index{precondition} of a function is a condition on its arguments that the caller must
make true before the call. A \emph{postcondition}\index{postcondition} is a condition on the result (and on the
arguments) that the function makes true when it returns, provided the precondition held.
\end{definition}

The two form a contract. If the caller keeps its side, the function must keep its own; if the caller breaks the
\omterm{def:g10:specifications-and-tests:pre-post}{precondition}, the function promises nothing at all.

\begin{omfigure}
\begin{tikzpicture}[font=\small, box/.style={draw, rounded corners=2pt, minimum height=9mm, inner sep=4pt}]
  \node[box, fill=omDef!8, text width=27mm, align=center] (caller) at (0,0) {caller};
  \node[box, fill=omThm!6, text width=27mm, align=center] (f) at (6.4,0) {function\\[-1pt]\footnotesize\code{leap_year}};
  \draw[ptr] (caller.north east) to[bend left=18] node[above, align=center, font=\footnotesize]
    {arguments that satisfy\\the \textbf{precondition}} (f.north west);
  \draw[ptr] (f.south west) to[bend left=18] node[below, align=center, font=\footnotesize]
    {a result that satisfies\\the \textbf{postcondition}} (caller.south east);
  \node[font=\footnotesize, align=left, anchor=west] at (8.3,0)
    {precondition broken:\\no promise at all};
\end{tikzpicture}

{\small The contract of a function. The caller's duty comes first, the function's duty after.}
\end{omfigure}

In Python the \omterm{def:g10:specifications-and-tests:specification}{specification} lives inside the function, as its first line of text.

\begin{definition}[Docstring]\label{def:g10:specifications-and-tests:docstring}
A \emph{docstring}\index{docstring} (documentation string) is a string written as the first statement of a
function's body. Python keeps it with the function, and \code{help(f)} prints it. In this book a docstring states
the function's \omterm{def:g10:specifications-and-tests:specification}{specification}.
\end{definition}

% ledger: g10:gregorian-rule, g10:gregorian-1582
The calendar in use almost everywhere today was instituted in $1582$. Its rule is that a year is a leap year when
it is divisible by $4$, except the century years ($1700$, $1800$, $1900$, \dots), which are leap years only when
they are divisible by $400$. Here is Sam's function, with a \omterm{def:g10:specifications-and-tests:docstring}{docstring} that says so.

\omcode{code/grade-10/01-specifications-and-tests/listings/leap.py}{1}{8}{Sam's function, with its
\omterm{def:g10:specifications-and-tests:specification}{specification}.}\label{lst:g10:specifications-and-tests:leap}

The \omterm{def:g10:specifications-and-tests:pre-post}{precondition} ($1583$ or later) is there because the rule means nothing before the calendar existed. The
\omterm{def:g10:specifications-and-tests:pre-post}{postcondition} says exactly which years must give \code{True}. With it, the run below is not a matter of opinion.

\omoutput{code/grade-10/01-specifications-and-tests/outputs/leap_demo.txt}{Four calls. The last line breaks the
\omterm{def:g10:specifications-and-tests:pre-post}{postcondition}: $1900$ is divisible by $100$ and not by $400$.}\label{out:g10:specifications-and-tests:leap-demo}

The body tests only the first part of the rule. The \omterm{def:g10:specifications-and-tests:specification}{specification} has turned ``my friend's calendar looks odd''
into a precise statement: for $1900$, the \omterm{def:g10:specifications-and-tests:pre-post}{postcondition} requires \code{False} and the function returns
\code{True}. That is what a bug is: a call whose \omterm{def:g10:specifications-and-tests:pre-post}{precondition} holds and whose \omterm{def:g10:specifications-and-tests:pre-post}{postcondition} fails.

\begin{method}[Writing a specification]\label{met:g10:specifications-and-tests:write-spec}
\begin{enumerate}
  \item Name each parameter and say what kind of value it holds (a whole number, a list of numbers, a string).
  \item Write the \omterm{def:g10:specifications-and-tests:pre-post}{precondition}: every restriction on the arguments (``not empty'', ``from 1 to 12'').
  \item Write the \omterm{def:g10:specifications-and-tests:pre-post}{postcondition} in terms of the arguments: what the result is, not how it is found.
  \item Check it on two or three examples worked by hand, including one at an edge.
\end{enumerate}
\end{method}

\begin{example}[The specification of a maximum]\label{ex:g10:specifications-and-tests:max-spec}
For a function \code{largest(t)}: the \omterm{def:g10:specifications-and-tests:pre-post}{precondition} is ``\code{t} is a non-empty list of numbers''; the
\omterm{def:g10:specifications-and-tests:pre-post}{postcondition} is ``the result is an element of \code{t}, and no element of \code{t} is larger''. The empty list
is excluded because it has no largest element: the \omterm{def:g10:specifications-and-tests:specification}{specification} forces the author to decide this case instead of
leaving it to chance.
\end{example}

The corrected version tests the exceptions first, from the rarest case to the commonest.

\omcode{code/grade-10/01-specifications-and-tests/listings/leap_fixed.py}{8}{12}{The body of the corrected
function (its \omterm{def:g10:specifications-and-tests:docstring}{docstring} is the same).}\label{lst:g10:specifications-and-tests:leap-fixed}

\section{Assertions}\label{sec:g10:specifications-and-tests:assertions}

A \omterm{def:g10:specifications-and-tests:specification}{specification} written in a \omterm{def:g10:specifications-and-tests:docstring}{docstring} is read by people. An \omterm{def:g10:specifications-and-tests:assertion}{assertion} is read by the machine.

\begin{definition}[Assertion]\label{def:g10:specifications-and-tests:assertion}
An \emph{assertion}\index{assertion} is a statement \code{assert condition, message} that evaluates the
condition. If it is true, the program goes on; if it is false, the program stops with an \code{AssertionError}
that carries the message.
\end{definition}

An \omterm{def:g10:specifications-and-tests:assertion}{assertion} at the top of a function checks its \omterm{def:g10:specifications-and-tests:pre-post}{precondition}: a caller who breaks the contract is stopped at once,
with a message that names the problem, instead of receiving a strange result much later.

\omcode{code/grade-10/01-specifications-and-tests/listings/average.py}{1}{14}{A \omterm{def:g10:specifications-and-tests:pre-post}{precondition} checked by an
\omterm{def:g10:specifications-and-tests:assertion}{assertion}. The last line breaks it on purpose.}\label{lst:g10:specifications-and-tests:average}

\omoutput{code/grade-10/01-specifications-and-tests/outputs/average.txt}{The first call keeps the
contract; the second stops with the message of the \omterm{def:g10:specifications-and-tests:assertion}{assertion}. The error report is read from the bottom: the
last line says what went wrong, the lines above say where.}\label{out:g10:specifications-and-tests:average}

Without the \omterm{def:g10:specifications-and-tests:assertion}{assertion}, \code{average([])} would stop too, but with \code{ZeroDivisionError} on the line
\code{total / len(t)}: a message about division, not about the caller's mistake.

\begin{remark}[What assertions are for]\label{rem:g10:specifications-and-tests:what-for}
An \omterm{def:g10:specifications-and-tests:assertion}{assertion} catches the \emph{programmer's} mistakes: a call that breaks a contract, a value that ``cannot''
happen. It is not the way to handle a user who types a wrong value: a program answers a user with a message and
asks again. One reason is that \omterm{def:g10:specifications-and-tests:assertion}{assertions} can be switched off.
% ledger: g10:py-assert-O
Python removes every \code{assert} statement when it is started with the option \code{-O}, so a program must
never rely on an \omterm{def:g10:specifications-and-tests:assertion}{assertion} to do its work.
\end{remark}

\omterm{def:g10:specifications-and-tests:assertion}{Assertions} often guard an expression that would crash: ``the list is not empty, \emph{and} its first element is
3''. This works because Python evaluates such a condition from left to right and stops as soon as the answer is
known.

\begin{definition}[Short-circuit evaluation]\label{def:g10:specifications-and-tests:short-circuit}
Python evaluates \code{a and b} by evaluating \code{a} first; if \code{a} is false, the result is false and
\code{b} is not evaluated at all. It evaluates \code{a or b} by evaluating \code{a} first; if \code{a} is true,
the result is true and \code{b} is not evaluated. This is \emph{short-circuit evaluation}\index{short-circuit
evaluation}.
\end{definition}

The function \code{loud} below prints a line each time it is evaluated, so the output shows which operands Python
actually computes.

\omcode{code/grade-10/01-specifications-and-tests/listings/short_circuit.py}{1}{11}{A function that announces
each evaluation.}\label{lst:g10:specifications-and-tests:short-circuit}

\omoutput{code/grade-10/01-specifications-and-tests/outputs/short_circuit.txt}{\code{loud} runs only in the middle
case, where the left operand does not settle the answer.}\label{out:g10:specifications-and-tests:short-circuit}

The order of the operands therefore matters. The guarded version below tests the length before it reads
\code{t[0]}; the unguarded version reads \code{t[0]} on an empty list.

\omcode{code/grade-10/01-specifications-and-tests/listings/guard.py}{1}{12}{A guard on the left of
\code{and}.}\label{lst:g10:specifications-and-tests:guard}

\omoutput{code/grade-10/01-specifications-and-tests/outputs/guard.txt}{The guarded function answers
\code{False} for the empty list; the unguarded one stops.}\label{out:g10:specifications-and-tests:guard}

Written \code{t[0] == x and len(t) > 0}, the guard would come too late: the crash happens while evaluating the
left operand.

\section{Choosing test cases}\label{sec:g10:specifications-and-tests:choosing}

Sam's three tries were not careless; they were badly chosen. None of them was a century year that is not a
multiple of $400$, such as $1900$, the only kind her function gets wrong. Choosing tests is a skill with a method.

\begin{definition}[Test case, boundary case]\label{def:g10:specifications-and-tests:test-case}
A \emph{test case}\index{test case} is an input for which the expected result is known in advance, from the
\omterm{def:g10:specifications-and-tests:specification}{specification}, together with that expected result. A test case \emph{passes} when the function returns the
expected result and \emph{fails} otherwise. A \emph{boundary case}\index{boundary case} is a test case at the edge
between two kinds of input that the \omterm{def:g10:specifications-and-tests:specification}{specification} treats differently.
\end{definition}

The \omterm{def:g10:specifications-and-tests:specification}{specification} of \code{leap_year} splits the years into four kinds, nested inside each other, as the figure
shows.

\begin{omfigure}
\begin{tikzpicture}[font=\small]
  \draw[thick, fill=black!3, rounded corners=6pt] (0,0) rectangle (11,4.4);
  \draw[thick, fill=omDef!8, rounded corners=6pt] (2.2,0.35) rectangle (10.6,4.05);
  \draw[thick, fill=omThm!8, rounded corners=6pt] (5.0,0.7) rectangle (10.25,3.7);
  \draw[thick, fill=omProp!14, rounded corners=6pt] (7.6,1.05) rectangle (9.9,3.35);
  \node[anchor=north west, align=left] at (0.15,4.3) {all years\\from 1583};
  \node[anchor=north west, align=left] at (2.35,3.95) {divisible\\by 4};
  \node[anchor=north west, align=left] at (5.15,3.6) {divisible\\by 100};
  \node[anchor=north west, align=left] at (7.75,3.25) {divisible\\by 400};
  \node[font=\footnotesize\ttfamily] at (1.1,1.2) {2023};
  \node[font=\footnotesize] at (1.1,0.75) {not leap};
  \node[font=\footnotesize\ttfamily] at (3.6,1.6) {2024};
  \node[font=\footnotesize] at (3.6,1.15) {leap};
  \node[font=\footnotesize\ttfamily] at (6.3,1.9) {1900};
  \node[font=\footnotesize] at (6.3,1.45) {not leap};
  \node[font=\footnotesize\ttfamily] at (8.75,2.0) {2000};
  \node[font=\footnotesize] at (8.75,1.55) {leap};
\end{tikzpicture}

{\small The four kinds of year in the \omterm{def:g10:specifications-and-tests:specification}{specification} of \code{leap_year}, with one test year in each. Sam's
three tries came from the two outer kinds and from the innermost one, never from the third.}
\end{omfigure}

\begin{method}[Choosing test cases from a specification]\label{met:g10:specifications-and-tests:choose}
\begin{enumerate}
  \item Split the allowed inputs into the kinds that the \omterm{def:g10:specifications-and-tests:specification}{specification} treats differently.
  \item Take one ordinary input of each kind.
  \item Add the \omterm{def:g10:specifications-and-tests:test-case}{boundary cases}: the inputs where one kind meets the next (a mark of exactly $10$ when $10$ or more
    passes), and the edges of the \omterm{def:g10:specifications-and-tests:pre-post}{precondition} (the smallest allowed year, a list of one element).
  \item If the \omterm{def:g10:specifications-and-tests:specification}{specification} says what happens to invalid inputs (an error, a special value), add one invalid
    input; if it says nothing, invalid inputs are not \omterm{def:g10:specifications-and-tests:test-case}{test cases}.
  \item Write the expected result of each case from the \omterm{def:g10:specifications-and-tests:specification}{specification}, \emph{before} running the function.
\end{enumerate}
\end{method}

\begin{example}[Tests for \code{leap_year}]\label{ex:g10:specifications-and-tests:leap-tests}
One year of each kind: $2023$, $2024$, $1900$, $2000$; the edge of the \omterm{def:g10:specifications-and-tests:pre-post}{precondition}: $1583$, which is not a leap
year. Sam's original function fails exactly one of these five tests, the one for $1900$. Any set of tests that misses the
third kind misses her bug, however many years it contains.
\end{example}

A set of tests should also make every line of the function run at least once. The corrected \code{leap_year} has
three \code{return} lines; the decision diagram below shows which test reaches which.

\begin{omfigure}
\begin{tikzpicture}[font=\small, q/.style={draw, rounded corners=2pt, fill=omDef!8, inner sep=3pt},
    r/.style={draw, fill=omProp!12, inner sep=3pt}, level distance=13mm]
  \node[q] (a) at (0,0) {\texttt{year \% 400 == 0}?};
  \node[r] (t1) at (-3,-1.4) {\code{return True}};
  \node[q] (b) at (2.4,-1.4) {\texttt{year \% 100 == 0}?};
  \node[r] (f1) at (-0.3,-2.8) {\code{return False}};
  \node[r] (l) at (4.8,-2.8) {\texttt{return year \% 4 == 0}};
  \draw[->] (a) -- node[above left, font=\footnotesize] {yes} (t1);
  \draw[->] (a) -- node[above right, font=\footnotesize] {no} (b);
  \draw[->] (b) -- node[above left, font=\footnotesize] {yes} (f1);
  \draw[->] (b) -- node[above right, font=\footnotesize] {no} (l);
  \node[below=1mm of t1, font=\footnotesize\ttfamily] {2000};
  \node[below=1mm of f1, font=\footnotesize\ttfamily] {1900};
  \node[below=1mm of l, font=\footnotesize\ttfamily] {2023, 2024, 1583};
\end{tikzpicture}

{\small The paths through the corrected function. The five tests reach every \code{return}; the last one in both
of its outcomes.}
\end{omfigure}

However well chosen, a set of tests tries finitely many inputs.

\begin{proposition}[Tests do not prove correctness]\label{prop:g10:specifications-and-tests:no-proof}
Let \code{f} be a function whose \omterm{def:g10:specifications-and-tests:pre-post}{precondition} allows infinitely many inputs, and let $T$ be any finite set of
\omterm{def:g10:specifications-and-tests:test-case}{test cases} that \code{f} passes. There is a function that also passes every test of $T$ and breaks the
\omterm{def:g10:specifications-and-tests:specification}{specification}.
\end{proposition}

\begin{proof}
Choose an allowed input $x$ that is not in $T$ (there is one, since $T$ is finite). Let \code{g} return what
\code{f} returns, except on $x$, where it returns a value that the \omterm{def:g10:specifications-and-tests:pre-post}{postcondition} forbids. On every input of $T$,
\code{g} returns what \code{f} returns, so it passes $T$; on $x$ it breaks the \omterm{def:g10:specifications-and-tests:specification}{specification}.
\end{proof}

Tests find bugs; they cannot show that none is left. Showing that a function is correct for \emph{every} input
needs a proof, which is the subject of a chapter next year. Tests remain the everyday tool: cheap, automatic, and
very good at catching the mistakes that people actually make.

\section{Tests that run themselves}\label{sec:g10:specifications-and-tests:automatic}

Tests typed by hand at the prompt are lost when the window closes. Written in a file, they can be run again after
every change, in a second, by anyone.

\begin{definition}[Test suite]\label{def:g10:specifications-and-tests:test-suite}
A \emph{test suite}\index{test suite} is a program that runs a set of \omterm{def:g10:specifications-and-tests:test-case}{test cases} on a function, compares each
result with the expected one, and reports which tests pass and which fail.
\end{definition}

Python's standard library includes \code{unittest} for this. Each test is a function whose name starts with
\code{test}, inside a class (the word \code{class} is explained in grade 12; here it only groups the tests).
\code{assertTrue(x)} checks that \code{x} is true, \code{assertFalse(x)} that it is false, and
\code{assertEqual(a, b)} that \code{a == b}. The first lines choose which version of \code{leap_year} to test.

\omcode{code/grade-10/01-specifications-and-tests/listings/leap_tests.py}{1}{27}{The five \omterm{def:g10:specifications-and-tests:test-case}{test cases} of the
example, as a \omterm{def:g10:specifications-and-tests:test-suite}{test suite}.}\label{lst:g10:specifications-and-tests:suite}

The session below runs the suite on Sam's version, then on the corrected one. (The runner also prints a line
with the time the tests took, which changes from run to run; \code{grep -v} removes it, so that the printed output
can be checked exactly. Commands like this one are explained next year.)

\omoutput{code/grade-10/01-specifications-and-tests/outputs/run_tests.txt}{A failing test is a bug report: it
names the test, the line, and what came back instead of what was
expected.}\label{out:g10:specifications-and-tests:run-tests}

A smaller tool, \code{doctest}, reads examples written in the \omterm{def:g10:specifications-and-tests:docstring}{docstring} itself, in the form of a session at the
prompt (\code{>>>} then the expected answer), and checks them.

\omcode{code/grade-10/01-specifications-and-tests/listings/mean.py}{1}{11}{Examples in a \omterm{def:g10:specifications-and-tests:docstring}{docstring} are also
tests.}\label{lst:g10:specifications-and-tests:mean}

\omoutput{code/grade-10/01-specifications-and-tests/outputs/doctest_run.txt}{\code{python -m doctest -v} runs each
example and compares what is printed.}\label{out:g10:specifications-and-tests:doctest}

\begin{inpractice}[Tests in a software project]
Professional projects keep thousands of tests and run all of them on every change, often with a tool such as
\code{pytest}, which finds the test functions by their names. A change is accepted only when every test passes.
When a bug is reported, the first step is to write a test that fails because of it; the bug is fixed when that
test passes, and the test stays in the suite so that the bug cannot come back unnoticed.
\end{inpractice}

\section{Where bugs come from}\label{sec:g10:specifications-and-tests:causes}

Most bugs belong to a handful of families. Knowing them tells you where to look first. One family deserves a
name.

\begin{definition}[Side effect]\label{def:g10:specifications-and-tests:side-effect}
A \emph{side effect}\index{side effect} of a function is any change it makes that outlives the call, apart from
the value it returns: a change to a list it received as an argument, to a variable outside the function, to a
file, or text printed on the screen.
\end{definition}

A \omterm{def:g10:specifications-and-tests:side-effect}{side effect} is not a mistake in itself (\code{print} exists for its \omterm{def:g10:specifications-and-tests:side-effect}{side effect}). It becomes a bug when the
\omterm{def:g10:specifications-and-tests:specification}{specification} does not mention it and the caller does not expect it.

\omcode{code/grade-10/01-specifications-and-tests/listings/side_effect.py}{1}{9}{The \omterm{def:g10:specifications-and-tests:specification}{specification} promises a
result, not a change to \code{t}.}\label{lst:g10:specifications-and-tests:side-effect}

\omoutput{code/grade-10/01-specifications-and-tests/outputs/side_effect.txt}{The result is right, but the
caller's list of scores has been reordered.}\label{out:g10:specifications-and-tests:side-effect}

The method \code{sort} changes the list it is applied to, and \code{t} is the caller's list itself, not a copy
(the next chapter, \ref{ch:g10:tuples-sets-comprehensions}, explains why). The fix is \code{sorted(t, reverse=True)[:3]}, which builds a new list.

A second family is the \code{if} that forgets a case.

\omcode{code/grade-10/01-specifications-and-tests/listings/missing_branch.py}{1}{12}{Three branches, and a mark
that falls between them.}\label{lst:g10:specifications-and-tests:missing-branch}

\omoutput{code/grade-10/01-specifications-and-tests/outputs/missing_branch.txt}{A function that reaches its end
without \code{return} returns \code{None}.}\label{out:g10:specifications-and-tests:missing-branch}

The \omterm{def:g10:specifications-and-tests:test-case}{boundary case} $15$ falls through all three tests. Writing the last branch as \code{else} makes the
\code{if} exhaustive, and the boundary test of the method above would have caught it.

\begin{method}[A checklist of common causes]\label{met:g10:specifications-and-tests:checklist}
When a test fails, or before trusting a function, look for:
\begin{enumerate}
  \item \emph{a value of the wrong type}: a number read with \code{input} is a string, and \code{"10" < "9"} is
    true because strings compare letter by letter;
  \item \emph{an unexpected \omterm{def:g10:specifications-and-tests:side-effect}{side effect}}: a list argument changed in place;
  \item \emph{an index out of range}: \code{t[len(t)]}, \code{t[0]} on an empty list, a loop one step too long;
  \item \emph{a non-exhaustive \code{if}}: a case that no branch handles;
  \item \emph{the wrong inequality}: \code{<} where the \omterm{def:g10:specifications-and-tests:specification}{specification} says ``at most'' (\code{<=});
  \item \emph{equality between decimal numbers}: computed decimals are rarely exact, so \code{==} between two of
    them is a trap (chapter \ref{ch:g10:floating-point} explains why);
  \item \emph{a misleading name}: a variable called \code{total} that holds an average invites the next mistake.
\end{enumerate}
\end{method}

\section{Exercises}\label{sec:g10:specifications-and-tests:exercises}

\begin{exercise}[$\star$]\label{exo:g10:specifications-and-tests:1}
The \omterm{def:g10:specifications-and-tests:docstring}{docstring} of a function \code{half(n)} holds the two examples \code{>>> half(10)} with the expected answer
\code{5} and \code{>>> half(7)} with the expected answer \code{3}. The body is \code{return n / 2}. What does
\code{python -m doctest} report for each example, and why?
\end{exercise}

\begin{exercise}[$\star$]\label{exo:g10:specifications-and-tests:2}
With the function \code{loud} of the listing on \omterm{def:g10:specifications-and-tests:short-circuit}{short-circuit evaluation}, what do these two lines print?
\code{print(3 > 5 and loud(1))} and \code{print(3 < 5 or loud(2))}. Then the same for
\code{print(loud(0) or loud(4))}.
\end{exercise}

\begin{exercise}[$\star$]\label{exo:g10:specifications-and-tests:3}
Write the \omterm{def:g10:specifications-and-tests:docstring}{docstring} of a function \code{smallest_positive(t)} that returns the smallest strictly positive number of
a list \code{t}. State a \omterm{def:g10:specifications-and-tests:pre-post}{precondition} and a \omterm{def:g10:specifications-and-tests:pre-post}{postcondition}, and decide what the \omterm{def:g10:specifications-and-tests:pre-post}{precondition} must say for the
function to always have an answer.
\end{exercise}

\begin{exercise}[$\star$]\label{exo:g10:specifications-and-tests:4}
The function \code{clamp(x, lo, hi)} has the \omterm{def:g10:specifications-and-tests:pre-post}{precondition} \code{lo <= hi} and returns \code{lo} if
\code{x < lo}, \code{hi} if \code{x > hi}, and \code{x} otherwise. List the kinds of input its \omterm{def:g10:specifications-and-tests:specification}{specification} treats
differently and the \omterm{def:g10:specifications-and-tests:test-case}{boundary cases} between them, and give one \omterm{def:g10:specifications-and-tests:test-case}{test case} (input and expected result) for each.
\end{exercise}

\begin{exercise}[$\star$]\label{exo:g10:specifications-and-tests:5}
The \omterm{def:g10:specifications-and-tests:specification}{specification} of \code{is_adult(age)} says: ``return \code{True} when \code{age} is $18$ or more''. The body is
\code{return age > 18}. Which \omterm{def:g10:specifications-and-tests:test-case}{test case} reveals the bug? Correct the body.
\end{exercise}

\begin{exercise}[$\star\star$]\label{exo:g10:specifications-and-tests:6}
The function below should return the non-negative numbers of a list.
\code{def remove_negatives(t):} then \code{for x in t:} \code{if x < 0: t.remove(x)} and finally
\code{return t}. Called on \code{[-1, -2, 3]} it returns \code{[-2, 3]}. Explain the result, name the two bugs (one
is a \omterm{def:g10:specifications-and-tests:side-effect}{side effect}), and write a version that builds a new list.
\end{exercise}

\begin{exercise}[$\star\star$]\label{exo:g10:specifications-and-tests:7}
The function \code{days_in_month(month, year)} returns the number of days of a month (\code{1} to \code{12}) of a
year from $1583$ on. Write a \code{unittest} \omterm{def:g10:specifications-and-tests:test-suite}{test suite} for it with every kind of input and every \omterm{def:g10:specifications-and-tests:test-case}{boundary case} you
can find, then a body that passes it (you may call the corrected \code{leap_year}).
\end{exercise}

\begin{exercise}[$\star\star$]\label{exo:g10:specifications-and-tests:8}
A function \code{is_sorted(t)} should return \code{True} when the list \code{t} is in increasing order. Its body is
\code{return len(t) < 2 or t[0] <= t[-1]}. For each \omterm{def:g10:specifications-and-tests:test-case}{test case} below, say whether it passes or fails:
\code{[]}, \code{[4]}, \code{[1, 2, 3]}, \code{[3, 2, 1]}, \code{[1, 5, 2, 6]}. Which kind of input did the author
forget?
\end{exercise}

\begin{exercise}[$\star\star$]\label{exo:g10:specifications-and-tests:9}
A friend says: ``my function \code{is_prime(n)} passes my tests for every \code{n} from $2$ to $1000$, so it is
correct''. Using the proposition of this chapter, build a function that passes the same tests and is wrong. What
would it take to know that \code{is_prime} is correct?
\end{exercise}

\begin{exercise}[$\star\star$]\label{exo:g10:specifications-and-tests:10}
The function \code{sign(x)} should return \code{-1}, \code{0} or \code{1} according to the sign of the number
\code{x}. Its body is \code{if x > 0: return 1} followed by \code{elif x < 0: return -1}. What does
\code{sign(0)} return? Which item of the checklist is this? Correct it.
\end{exercise}

\begin{exercise}[$\star\star$]\label{exo:g10:specifications-and-tests:11}
Write a function \code{percent(part, whole)} that returns \code{part} as a percentage of \code{whole}. Its
\omterm{def:g10:specifications-and-tests:pre-post}{preconditions} are \code{whole > 0} and \code{0 <= part <= whole}; check them with \omterm{def:g10:specifications-and-tests:assertion}{assertions} with clear messages, and
check the \omterm{def:g10:specifications-and-tests:pre-post}{postcondition} (the result lies between $0$ and $100$) with a third \omterm{def:g10:specifications-and-tests:assertion}{assertion}.
\end{exercise}

\begin{exercise}[$\star\star\star$]\label{exo:g10:specifications-and-tests:12}
A ticket costs $10$; a child under $12$ pays half, and on Tuesdays everyone has $2$ off (after the child
reduction). The function has two \code{if} statements one after the other, one for each rule. \emph{(a)} How many
different paths through the function are there? \emph{(b)} What is the smallest number of tests that makes every
line run? \emph{(c)} What is the smallest number that follows every path? Give such tests with their expected
prices.
\end{exercise}

\begin{exercise}[$\star\star\star$]\label{exo:g10:specifications-and-tests:13}
Write \code{roman(n)} for $1 \le n \le 3999$, which returns the Roman numeral of \code{n} (\code{4} is \code{IV},
\code{9} is \code{IX}, \code{1994} is \code{MCMXCIV}), with its \omterm{def:g10:specifications-and-tests:specification}{specification}, an \omterm{def:g10:specifications-and-tests:assertion}{assertion} on its \omterm{def:g10:specifications-and-tests:pre-post}{precondition} and a
\omterm{def:g10:specifications-and-tests:test-suite}{test suite} that includes the subtractive forms, the smallest and the largest allowed values, and two invalid inputs.
\end{exercise}

\begin{exercise}[$\star\star\star$]\label{exo:g10:specifications-and-tests:14}
This loop should count the steps of $0.1$ from $0$ to $1$: \code{x = 0.0}, \code{steps = 0}, then
\code{while x != 1.0:} with the body \code{x = x + 0.1} and \code{steps = steps + 1}. Run in your head with exact
decimals, it stops after $10$ steps; run by Python, it never stops. Which item of the checklist explains it? Fix the
loop in two ways: with a counter that does not depend on \code{x}, and with \code{math.isclose(x, 1.0)} as the test.
\end{exercise}

\begin{exercise}[$\star\star\star$]\label{exo:g10:specifications-and-tests:15}
A \emph{mutant} of a function is a copy with one small change, such as one number or one operator replaced. A \omterm{def:g10:specifications-and-tests:test-suite}{test
suite} \emph{kills} a mutant when at least one of its tests fails on it. For the corrected \code{leap_year} and the
suite of the two years $2024$ and $2000$, say which of these mutants survive: \emph{(a)} \code{400} replaced by
\code{100}; \emph{(b)} \code{return False} replaced by \code{return True}; \emph{(c)} the last line's \code{4}
replaced by \code{400}; \emph{(d)} \texttt{year \% 100 == 0} replaced by \texttt{year \% 100 != 0}. Which single year
added to the suite kills every survivor?
\end{exercise}

\section{Problem: The Mutants of the Leap Year}\label{sec:g10:specifications-and-tests:problem}

\begin{problem}[{Weekend problem --- how good is a test suite?}]\label{pb:g10:specifications-and-tests:1}
Sam's teacher grades \omterm{def:g10:specifications-and-tests:test-suite}{test suites}, not programs. Her method: she makes many mutants of a correct function (see
\cref{exo:g10:specifications-and-tests:15}) and counts how many of them a suite kills. A suite that kills every mutant
would have caught all those mistakes. This problem compares two suites for the corrected \code{leap_year}: four
years chosen with the method of this chapter, and a hundred years drawn at random between $1583$ and $2400$.

\needspace{4\baselineskip}\medskip\noindent\textbf{Part I --- The specification.}
\begin{enumerate}
  \item State the \omterm{def:g10:specifications-and-tests:pre-post}{precondition} of \code{leap_year}, and explain why it starts at $1583$.
  \item Name the four kinds of year that the \omterm{def:g10:specifications-and-tests:pre-post}{postcondition} treats differently, and say for each whether it is a
    leap year.
  \item How many years from $1583$ to $2400$ (both included) are century years? How many of them are leap years?
  \item What fraction of the years from $1583$ to $2400$ belong to the kind that Sam's first version got wrong?
\end{enumerate}

\needspace{4\baselineskip}\medskip\noindent\textbf{Part II --- A suite from the specification.}
\begin{enumerate}[resume]
  \item Give the smallest suite with one year of each kind, using $2023$, $2024$, $1900$ and $2000$, with the
    expected results.
  \item Show that no suite of three years can contain one year of each kind.
  \item Which test of this suite does Sam's first version fail, and which of the four kinds would a suite need to
    miss to let her bug through?
  \item Why is $1583$ a \omterm{def:g10:specifications-and-tests:test-case}{boundary case}, and what is the expected result for it?
\end{enumerate}

\needspace{4\baselineskip}\medskip\noindent\textbf{Part III --- Mutants.}
\begin{enumerate}[resume]
  \item The teacher's program changes the body of the corrected function (its lines $8$ to $12$) in one place, by
    these rules: each \code{==} becomes \code{!=}; each \code{True} becomes \code{False} and each \code{False} becomes
    \code{True}; the number \code{400} becomes \code{100} or \code{4}; \code{100} becomes \code{400} or \code{4};
    \code{4} becomes \code{100} or \code{400}; each \code{0} on the right of an \code{==} becomes \code{1}. Count
    the mutants by hand.
  \item Write out the mutant where \texttt{year \% 100 == 0} becomes \texttt{year \% 100 != 0}, and find one year of the
    four-year suite on which it fails.
  \item Write out the mutant where \texttt{year \% 400} becomes \texttt{year \% 4}. Which kind of year does it get wrong?
  \item Run the program \code{exercises/mutants.py} of this chapter's code. How many mutants does the four-year suite
    kill?
  \item Explain without running it why the mutant ``the first \code{True} becomes \code{False}'' is killed by the
    four-year suite.
\end{enumerate}

\needspace{4\baselineskip}\medskip\noindent\textbf{Part IV --- A hundred random years.}
\begin{enumerate}[resume]
  \item Using the answer to question 3, how many century years would you expect among $100$ years drawn at random
    from $1583$ to $2400$, as a rough frequency?
  \item The random suite of the program kills only $7$ mutants. List the seven survivors that the program prints and
    find what they have in common.
  \item How many century years does the program's random suite contain? Explain the survivors with this number.
  \item Which single year added to the random suite would kill the most survivors? Does any one year kill all seven?
  \item Find the smallest set of years that, added to the random suite, kills all seven survivors.
  \item State the result of the comparison: how many mutants does each suite kill, and with how many tests?
\end{enumerate}
\end{problem}
